\documentclass[12pt]{article}
\usepackage{amsmath, amssymb, amsthm}
\usepackage{geometry}
\geometry{margin=1in}

\title{Math 4102: Introduction to Lebesgue Integration \\ Final Examination Problems List}
\author{Aaron Chen}
\date{May 2, 2026}

\begin{document}
\maketitle

\section*{Problem 1 (30.1, p.85)}
State and prove Egoroff's theorem.

\begin{proof}
\end{proof}

\section*{Problem 2 (30.3, p.87)}
State and prove Lusin's theorem.

\begin{proof}
\end{proof}

\section*{Problem 3}
State the definition of a normed linear space and give both finite- and $\infty$-dimensional examples.

\begin{proof}
\end{proof}

\section*{Problem 4}
State the definition of a real Hilbert space and give both finite- and $\infty$-dimensional examples.

\begin{proof}
\end{proof}

\section*{Problem 5}
Determine whether $\{\cos kx : k=0,1,2...\}$ is an orthogonal set in $\mathcal{L}^{2}[0,\pi]$.

\begin{proof}
\end{proof}

\section*{Problem 6 (32.13, p.98)}
State and prove Cauchy's inequality.

\begin{proof}
\end{proof}

\section*{Problem 7 (33.9, p.104)}
State and prove Minkowski's inequality in $\mathcal{L}^{2}$.

\begin{proof}
\end{proof}

\section*{Problem 8 (p.113)}
Define a trigonometric polynomial and a trigonometric series.

\begin{proof}
\end{proof}

\section*{Problem 9}
State and prove Bessel's inequality for an orthonormal set in $\mathcal{L}^{2}[-\pi,\pi]$.

\begin{proof}
\end{proof}

\section*{Problem 10}
State and prove the Riemann-Lebesgue lemma for the Fourier series of $f \in \mathcal{L}^{2}[-\pi,\pi]$.

\begin{proof}
\end{proof}

\section*{Problem 11 (37.4, p.121)}
Give two formulas for the Dirichlet kernel $D_{n}(t)$.

\begin{proof}
\end{proof}

\section*{Problem 12 (37.5, p.123)}
Express the partial sum $s_{n}(x)$ of the Fourier series for a function $f \in \mathcal{L}^{2}[-\pi,\pi]$ using the Lebesgue integral and the Dirichlet kernel.

\begin{proof}
\end{proof}

\section*{Problem 13 (37.6, p.124)}
Give a formula for the Fejér kernel $K_{n}(t)$ and relate it to the Dirichlet kernel.

\begin{proof}
\end{proof}

\section*{Problem 14}
Define the Cesàro $(C,1)$-sum of a series $\sum_{k=1}^{\infty} a_{k}$. Give an example of a divergent series with a convergent $(C, 1)$-sum.

\begin{proof}
\end{proof}

\section*{Problem 15 (37.6, p.124)}
Express the $(C, 1)$-sum $\sigma_{n}(x)$ of the Fourier series for a function $f \in \mathcal{L}[-\pi,\pi]$ using the Lebesgue integral and the Fejér kernel.

\begin{proof}
\end{proof}

\section*{Problem 16 (37.8, p.125)}
Prove that if $f \in \mathcal{L}^{2}[-\pi,\pi]$ and $f$ is continuous, then $\{\sigma_{n}\}$ converges uniformly to $f$.

\begin{proof}
\end{proof}

\section*{Problem 17 (40.2, p.139)}
Prove that if $f \in \mathcal{L}^{2}[-\pi,\pi]$ and $f$ is differentiable at $x$, then $\{s_{n}(x)\}$ converges to $f(x)$.

\begin{proof}
\end{proof}

\section*{Problem 18 (41.2, p.141)}
State and prove the Riemann localization lemma.

\begin{proof}
\end{proof}

\section*{Problem 19}
Compute the complex exponential Fourier series of $f(x) = \sin x + 2 \cos 3x$.

\begin{proof}
\end{proof}

\end{document}